\documentclass[10pt,a4paper]{amsart}
\usepackage{amsmath,amssymb,amsthm,graphicx}
\usepackage{verbatim}
\newenvironment{reference}{\comment}{\endcomment}
%\newenvironment{reference}{}{}
\newenvironment{slogan}{\comment}{\endcomment}
\newenvironment{history}{\comment}{\endcomment}

% For commutative diagrams we use Xy-pic
\usepackage[all]{xy}

% We use 2cell for 2-commutative diagrams.
\xyoption{2cell}
% Theorem environments.
%
\theoremstyle{plain}
\newtheorem{theorem}[subsection]{Theorem}
\newtheorem{proposition}[subsection]{Proposition}
\newtheorem{lemma}[subsection]{Lemma}

\theoremstyle{definition}
\newtheorem{definition}[subsection]{Definition}
\newtheorem{example}[subsection]{Example}
\newtheorem{exercise}[subsection]{Exercise}
\newtheorem{situation}[subsection]{Situation}

\theoremstyle{remark}
\newtheorem{remark}[subsection]{Remark}
\newtheorem{remarks}[subsection]{Remarks}

\numberwithin{equation}{subsection}

% Macros
%
\def\lim{\mathop{\mathrm{lim}}\nolimits}
\def\colim{\mathop{\mathrm{colim}}\nolimits}
\def\Spec{\mathop{\mathrm{Spec}}}
\def\Hom{\mathop{\mathrm{Hom}}\nolimits}
\def\Ext{\mathop{\mathrm{Ext}}\nolimits}
\def\SheafHom{\mathop{\mathcal{H}\!\mathit{om}}\nolimits}
\def\SheafExt{\mathop{\mathcal{E}\!\mathit{xt}}\nolimits}
\def\Sch{\mathit{Sch}}
\def\Mor{\mathop{\mathrm{Mor}}\nolimits}
\def\Ob{\mathop{\mathrm{Ob}}\nolimits}
\def\Sh{\mathop{\mathit{Sh}}\nolimits}
\def\NL{\mathop{N\!L}\nolimits}
\def\CH{\mathop{\mathrm{CH}}\nolimits}
\def\proetale{{pro\text{-}\acute{e}tale}}
\def\etale{{\acute{e}tale}}
\def\QCoh{\mathit{QCoh}}
\def\Ker{\mathop{\mathrm{Ker}}}
\def\Im{\mathop{\mathrm{Im}}}
\def\Coker{\mathop{\mathrm{Coker}}}
\def\Coim{\mathop{\mathrm{Coim}}}

% Boxtimes
%
\DeclareMathSymbol{\boxtimes}{\mathbin}{AMSa}{"02}

%
% Macros for moduli stacks/spaces
%
\def\QCohstack{\mathcal{QC}\!\mathit{oh}}
\def\Cohstack{\mathcal{C}\!\mathit{oh}}
\def\Spacesstack{\mathcal{S}\!\mathit{paces}}
\def\Quotfunctor{\mathrm{Quot}}
\def\Hilbfunctor{\mathrm{Hilb}}
\def\Curvesstack{\mathcal{C}\!\mathit{urves}}
\def\Polarizedstack{\mathcal{P}\!\mathit{olarized}}
\def\Complexesstack{\mathcal{C}\!\mathit{omplexes}}
% \Pic is the operator that assigns to X its picard group, usage \Pic(X)
% \Picardstack_{X/B} denotes the Picard stack of X over B
% \Picardfunctor_{X/B} denotes the Picard functor of X over B
\def\Pic{\mathop{\mathrm{Pic}}\nolimits}
\def\Picardstack{\mathcal{P}\!\mathit{ic}}
\def\Picardfunctor{\mathrm{Pic}}
\def\Deformationcategory{\mathcal{D}\!\mathit{ef}}
\usepackage{cleveref}
\crefname{originalreference}{original reference}{original references}
\Crefname{originalreference}{Original reference}{Original references}
\makeatletter
\newlabel{definition-rig-etale}{{0AQM}{source}{}{sourcepointer-1}{}}
\newlabel{definition-rig-etale@cref}{{[originalreference][0][]0AQM}{[1][0][]source}{}{}{}}
\makeatother
\crefname{originaldefinition}{original definition}{original definitions}
\Crefname{originaldefinition}{Original definition}{Original definitions}
\makeatletter
\newlabel{formal-spaces-definition-formal-algebraic-space}{{0AIM}{source}{}{sourcepointer-2}{}}
\newlabel{formal-spaces-definition-formal-algebraic-space@cref}{{[originaldefinition][0][]0AIM}{[1][0][]source}{}{}{}}
\makeatother
\crefname{originallemma}{original lemma}{original lemmas}
\Crefname{originallemma}{Original lemma}{Original lemmas}
\makeatletter
\newlabel{formal-spaces-lemma-property-defines-property-morphisms}{{0ANG}{source}{}{sourcepointer-3}{}}
\newlabel{formal-spaces-lemma-property-defines-property-morphisms@cref}{{[originallemma][0][]0ANG}{[1][0][]source}{}{}{}}
\makeatother
\crefname{originalremark}{original remark}{original remarks}
\Crefname{originalremark}{Original remark}{Original remarks}
\makeatletter
\newlabel{formal-spaces-remark-variant-Noetherian}{{0ANI}{source}{}{sourcepointer-4}{}}
\newlabel{formal-spaces-remark-variant-Noetherian@cref}{{[originalremark][0][]0ANI}{[1][0][]source}{}{}{}}
\makeatother
\crefname{originalreference}{original reference}{original references}
\Crefname{originalreference}{Original reference}{Original references}
\makeatletter
\newlabel{lemma-rig-etale-axioms}{{0AQL}{source}{}{sourcepointer-5}{}}
\newlabel{lemma-rig-etale-axioms@cref}{{[originalreference][0][]0AQL}{[1][0][]source}{}{}{}}
\makeatother
\title{Four equivalent characterizations of rig-etale morphisms}
\author{The Stacks Project Authors}
\date{stacks-project:restricted.tex:local-snapshot-20261012; GNU FDL 1.2 or later}
\begin{document}
\maketitle
\noindent\textit{Editorial note.} All delivered mathematical text is the author's own, copied from the pinned Stacks Project chapter file restricted.tex; no statement or proof step was rewritten, reordered or removed.\par\medskip
\noindent\textit{Editorial note.} The counted proof is the author's complete argument: the proof is the standard property-defines-property-morphisms template with locality of rig-etale-ness as the input property\par\medskip
\noindent\textit{Editorial note.} The label lemma-rig-etale-morphisms names the counted statement, and the proof environment immediately following it in the source is delivered verbatim as its proof.\par\medskip
\noindent\textit{Editorial note.} The counted statement is carried by the author's own lemma environment in the chapter source.\par\medskip
\noindent\textit{Allowed standard prerequisites (editorial).}\begin{itemize}
\item the definition of rig-etale morphisms of locally Noetherian formal algebraic spaces via etale covers and rig-etale arrows of WAdm\textasciicircum{}Noeth (definition-rig-etale)
\item rig-etale-ness is a local property of arrows of WAdm\textasciicircum{}Noeth (lemma-rig-etale-axioms of the same chapter)
\item the Formal Spaces lemma that a local property of arrows defines a property of morphisms (formal-spaces-lemma-property-defines-property-morphisms) in its Noetherian variant (formal-spaces-remark-variant-Noetherian)
\end{itemize}\medskip
\ifdefined\hypertarget\hypertarget{sourcepointer-1}{}\fi
\noindent\textit{Source reference.} Original reference 0AQM: sources/stacks-restricted/pinned-original-source.tex, source lines 5027--5027. Prerequisite pointer invoked at proof line(s) 5083; it names a standard fact of the same chapter that the counted argument uses, and every step of the counted proof built on it is retained verbatim in the delivered body.\par\medskip
\ifdefined\hypertarget\hypertarget{sourcepointer-2}{}\fi
\noindent\textit{Source reference.} Original definition 0AIM: sources/stacks-chow/tags.txt, source lines 13270--13270. Prerequisite pointer invoked at proof line(s) 5065; it names a standard fact of a sibling chapter that the counted argument uses, and every step of the counted proof built on it is retained verbatim in the delivered body.\par\medskip
\ifdefined\hypertarget\hypertarget{sourcepointer-3}{}\fi
\noindent\textit{Source reference.} Original lemma 0ANG: sources/stacks-chow/tags.txt, source lines 13439--13439. Prerequisite pointer invoked at proof line(s) 5089; it names a standard fact of a sibling chapter that the counted argument uses, and every step of the counted proof built on it is retained verbatim in the delivered body.\par\medskip
\ifdefined\hypertarget\hypertarget{sourcepointer-4}{}\fi
\noindent\textit{Source reference.} Original remark 0ANI: sources/stacks-chow/tags.txt, source lines 13441--13441. Prerequisite pointer invoked at proof line(s) 5092; it names a standard fact of a sibling chapter that the counted argument uses, and every step of the counted proof built on it is retained verbatim in the delivered body.\par\medskip
\ifdefined\hypertarget\hypertarget{sourcepointer-5}{}\fi
\noindent\textit{Source reference.} Original reference 0AQL: sources/stacks-restricted/pinned-original-source.tex, source lines 4788--4788. Prerequisite pointer invoked at proof line(s) 5086; it names a standard fact of the same chapter that the counted argument uses, and every step of the counted proof built on it is retained verbatim in the delivered body.\par\medskip
\begin{lemma}
\label{lemma-rig-etale-morphisms}
Let $S$ be a scheme. Let $f : X \to Y$ be a morphism of
locally Noetherian formal algebraic spaces over $S$.
The following are equivalent
\begin{enumerate}
\item $f$ is rig-\'etale,
\item for every commutative diagram
$$
\xymatrix{
U \ar[d] \ar[r] & V \ar[d] \\
X \ar[r] & Y
}
$$
with $U$ and $V$ affine formal algebraic spaces, $U \to X$ and $V \to Y$
representable by algebraic spaces and \'etale, the morphism $U \to V$
corresponds to a rig-\'etale map in $\textit{WAdm}^{Noeth}$,
\item there exists a covering $\{Y_j \to Y\}$ as in
Formal Spaces,
Definition \ref{formal-spaces-definition-formal-algebraic-space}
and for each $j$
a covering $\{X_{ji} \to Y_j \times_Y X\}$ as in
Formal Spaces,
Definition \ref{formal-spaces-definition-formal-algebraic-space}
such that each $X_{ji} \to Y_j$  corresponds
to a rig-\'etale map in $\textit{WAdm}^{Noeth}$, and
\item there exist a covering $\{X_i \to X\}$ as in
Formal Spaces,
Definition \ref{formal-spaces-definition-formal-algebraic-space}
and for each $i$ a factorization $X_i \to Y_i \to Y$ where $Y_i$
is an affine formal algebraic space, $Y_i \to Y$ is representable
by algebraic spaces and \'etale, and $X_i \to Y_i$ corresponds
to a rig-\'etale map in $\textit{WAdm}^{Noeth}$.
\end{enumerate}
\end{lemma}

\begin{proof}
The equivalence of (1) and (2) is Definition \ref{definition-rig-etale}.
The equivalence of (2), (3), and (4) follows from the fact that
being rig-\'etale is a local property of arrows of
$\text{WAdm}^{Noeth}$ by Lemma \ref{lemma-rig-etale-axioms}
and an application of the variant of
Formal Spaces, Lemma
\ref{formal-spaces-lemma-property-defines-property-morphisms}
for morphisms between locally Noetherian algebraic spaces
mentioned in Formal Spaces, Remark
\ref{formal-spaces-remark-variant-Noetherian}.
\end{proof}

\end{document}
